\chapter{Integrability And The Rigid Body}
\label{chap:rigid-integrability}

\section{The Rigid Body As A Bridge To Integrability}
\label{sec:rigid-why-central}

The free rigid body is a compact model in which several central ideas meet. Its
motion is nonlinear; its configuration space is the Lie group \(\SO(3)\); its
reduced description is Hamiltonian on angular-momentum space; and it is
completely integrable. At the same time, small modifications such as gravity
acting on a heavy top lead toward nonintegrability and chaos. The classical
integrable exceptions beyond the free body are Lagrange's symmetric heavy top
(1788) and Kovalevskaya's top (1888), with other special cases requiring
additional hypotheses. Husson's 1906 analysis made this rarity precise for the
classical heavy top under its stated assumptions. Their rarity is part of the
lesson.

The most vivid symptom is the intermediate-axis flip. Spin a book about its
longest or shortest principal axis and the rotation is stable; spin it about
the middle axis and the book turns over. The behavior comes from the geometry
of two conserved quantities, energy and angular momentum, intersecting on the
angular-momentum sphere. The rigid body is a place
where geometric pictures become part of the argument: reduction, stability,
Poinsot rolling, and invariant tori are all different views of the same motion.

\begin{assumption}[Rigid body model]
A rigid body is an ideal continuum whose internal distances do not change. A
material point is labelled by a body coordinate \(X\in B\subset\R^3\). Its
spatial position is
\[
  x(t)=r(t)+R(t)X,
\]
where \(r(t)\in\R^3\) is the center-of-mass position and
\(R(t)\in\SO(3)\) is the attitude matrix. The free rigid body chapter suppresses
translation by working in the center-of-mass frame.
\end{assumption}

\paragraph{Three layers to keep separate.}
Rigid-body mechanics mixes three descriptions that are often compressed into
one line of notation. The kinematic layer says how \(R(t)\in\SO(3)\) defines
angular velocity. The dynamical layer says how the kinetic energy and the
inertia tensor produce Euler's equations. The reduced Hamiltonian layer says
why angular momentum spheres carry a Lie-Poisson flow. Each layer has its own
variables and its own conserved quantities. Keeping the layers separate
preserves the distinction between body components and spatial components.

\paragraph{Roadmap.}
The chapter first derives the free rigid body's kinematics, Euler equations,
reduction, and stability, then abstracts the lesson into Liouville integrability
and invariant tori.

\section{The Rotation Group And Angular Velocity}
\label{sec:rotation-group-angular-velocity}

The rotation group is
\[
  \SO(3)=\{R\in\R^{3\times 3}:R^TR=I,\ \det R=1\}.
\]
The matrix \(R\) has columns equal to the three orthonormal body axes
\(\hat e_1,\hat e_2,\hat e_3\) as seen in space:
\[
  R=(\hat e_1\ \hat e_2\ \hat e_3).
\]
Thus a point of configuration space is not a triple of unconstrained vectors;
it is an orthonormal oriented frame. The condition \(R^TR=I\) gives six
constraints on nine matrix entries, leaving three degrees of freedom.

The spatial angular velocity \(\omega_{\mathrm{space}}\) is defined by the fact that every body
axis instantaneously rotates according to
\[
  \frac{\dd \hat e_a}{\dd t}=\omega_{\mathrm{space}}\times\hat e_a,
\]
or, in components,
\[
  \dot R_{ia}=\epsilon_{ijk}(\omega_{\mathrm{space}})_jR_{ka}.
\]

Differentiate \(R^TR=I\):
\[
  \dot R^TR+R^T\dot R=0.
\]
Thus \(R^T\dot R\) is skew-symmetric. Every skew-symmetric \(3\times3\) matrix
is uniquely of the form
\[
  \widehat\Omega=
  \begin{pmatrix}
  0 & -\Omega_3 & \Omega_2\\
  \Omega_3 & 0 & -\Omega_1\\
  -\Omega_2 & \Omega_1 & 0
  \end{pmatrix},
\]
where \(\widehat\Omega v=\Omega\times v\). The body angular velocity is defined
by
\[
  R^T\dot R=\widehat\Omega.
\]
The hat map is a Lie algebra isomorphism:
\[
  [\widehat a,\widehat b]=\widehat{a\times b}.
\]
This identity is the bridge between matrix commutators in \(\mathfrak{so}(3)\)
and vector cross products in the rigid-body equations.
The spatial angular velocity \(\omega_{\mathrm{space}}\) satisfies
\[
  \dot R R^T=\widehat{\omega}_{\mathrm{space}},
  \qquad
  \omega_{\mathrm{space}}=R\Omega.
\]

This space/body distinction is one of the main sources of sign errors in rigid
body mechanics. The vector \(\omega_{\mathrm{space}}\) is expressed in fixed
space axes, while \(\Omega\) is expressed in the moving body axes. They describe
the same instantaneous rotation, but derivatives of body components must also
differentiate the moving frame. Euler's equation is where this bookkeeping
becomes dynamical.

\paragraph{Derivative rules for moving frames.}
If \(A(t)\) is a vector written in body components, then its spatial vector is
\(a(t)=R(t)A(t)\). Differentiating gives
\[
  \dot a
  =
  \dot R A+R\dot A
  =
  R(\widehat\Omega A+\dot A)
  =
  R(\dot A+\Omega\times A).
\]
Thus
\[
  \boxed{
  \dot a_{\mathrm{space}}
  =
  R\left(\dot A_{\mathrm{body}}+\Omega\times A_{\mathrm{body}}\right).
  }
\]
Conversely, if \(a\) is fixed in space and \(A=R^Ta\) are its body components,
then
\[
  \dot A
  =
  -\widehat\Omega A
  =
  -\Omega\times A.
\]
This rule explains the form of Euler's equation before any variational
calculation is done: a vector can be conserved in space while its body
components precess because the axes themselves are moving.

\section{\texorpdfstring{Euler Angles, Metric, And The Topology Of \(\SO(3)\)}{Euler Angles, Metric, And The Topology Of SO(3)}}
\label{sec:euler-angles-so3}

One can parametrize a dense open set of \(\SO(3)\) by Euler angles
\[
  (\phi,\theta,\psi),
  \qquad
  0\leq \phi<2\pi,\quad 0\leq \theta\leq \pi,\quad 0\leq \psi<2\pi.
\]
For the Euler-angle convention used below, the body axes are
\[
\begin{aligned}
\hat e_1 &=
\begin{pmatrix}
\cos\phi\cos\psi-\sin\phi\cos\theta\sin\psi\\
\sin\phi\cos\psi+\cos\phi\cos\theta\sin\psi\\
\sin\theta\sin\psi
\end{pmatrix},\\
\hat e_2 &=
\begin{pmatrix}
-\cos\phi\sin\psi-\sin\phi\cos\theta\cos\psi\\
-\sin\phi\sin\psi+\cos\phi\cos\theta\cos\psi\\
\sin\theta\cos\psi
\end{pmatrix},\\
\hat e_3 &=
\begin{pmatrix}
\sin\phi\sin\theta\\
-\cos\phi\sin\theta\\
\cos\theta
\end{pmatrix}.
\end{aligned}
\]
The convention matters for signs in formulas, but the lesson is geometric:
Euler angles are a coordinate chart, not the space itself. The chart is singular at
\(\theta=0,\pi\), even though \(\SO(3)\) is a smooth homogeneous space.

A coordinate-independent line element on \(\SO(3)\) is
\[
  \dd s^2=-\frac12\tr\left(R^{-1}\dd R\right)^2.
\]
This is the bi-invariant metric induced by minus the Killing-form normalization
on \(\mathfrak{so}(3)\); the minus sign makes the square norm positive because
skew matrices have negative trace square.
For this standard Euler-angle convention, this becomes
\[
  \dd s^2
  =
  \dd\phi^2+\dd\theta^2+\dd\psi^2
  +2\cos\theta\,\dd\phi\,\dd\psi.
\]
Introduce
\[
  \xi_1=\phi+\psi,\qquad \xi_2=\phi-\psi.
\]
Then
\[
  \dd s^2
  =
  \dd\theta^2
  +
  \cos^2\frac{\theta}{2}\,\dd\xi_1^2
  +
  \sin^2\frac{\theta}{2}\,\dd\xi_2^2.
\]
This formula shows the topology pictorially. At one endpoint of the
\(\theta\)-interval the \(\xi_2\)-circle shrinks; at the other endpoint the
\(\xi_1\)-circle shrinks. Gluing these solid tori gives \(S^3\), and the
remaining identification is the antipodal identification. Hence
\[
  \SO(3)\simeq S^3/\Z_2.
\]
This is why a rigid body admits local angle coordinates but no single global
nondegenerate Euler-angle chart.

\section{Kinetic Energy And The Inertia Tensor}
\label{sec:inertia-tensor}

Take the center of mass as origin in body coordinates:
\[
  \int_B X\,\rho(X)\,\dd^3X=0,
\]
where \(\rho(X)\) is the mass density. A point \(X\) has velocity
\[
  \dot x=\dot R X=R\widehat\Omega X=R(\Omega\times X).
\]
Since \(R\) preserves Euclidean lengths,
\[
  \norm{\dot x}^2=\norm{\Omega\times X}^2.
\]
The kinetic energy is
\[
  T=\frac12\int_B\rho(X)\norm{\Omega\times X}^2\,\dd^3X.
\]
Use the vector identity
\[
  \norm{\Omega\times X}^2
  =
  \norm{\Omega}^2\norm{X}^2-(\Omega\cdot X)^2,
\]
to write the kinetic energy as
\[
  T=\frac12\Omega^T I\Omega,
\]
where
\[
  I_{ij}
  =
  \int_B\rho(X)
  \left(\norm{X}^2\delta_{ij}-X_iX_j\right)\dd^3X.
\]
The matrix \(I\) is the inertia tensor in the body frame. In principal axes,
\[
  I=\diag(I_1,I_2,I_3),
\]
and
\[
  T=\frac12(I_1\Omega_1^2+I_2\Omega_2^2+I_3\Omega_3^2).
\]
Equivalently, using the spatial angular velocity \(\omega_{\mathrm{space}}\) and the body axes
\(\hat e_a\),
\[
  L=\frac12\sum_{a=1}^3 I_a(\hat e_a\cdot\omega_{\mathrm{space}})^2.
\]
Here \(L\) denotes the Lagrangian. Later symbols such as \(L_z\) denote angular
momentum components; the subscripted notation is always specified locally.
The spatial angular momentum is
\[
  J_i=\frac{\partial L}{\partial (\omega_{\mathrm{space}})_i}
  =
  \sum_{a=1}^3 I_a(\hat e_a)_i(\hat e_a\cdot\omega_{\mathrm{space}}).
\]
Its body components are
\[
  \widetilde J_a=\hat e_a\cdot J=I_a\Omega_a,
\]
and the Hamiltonian is
\[
  H=\sum_{a=1}^3\frac{\widetilde J_a^2}{2I_a}.
\]

\section{\texorpdfstring{Full Phase Space: \(T^*\SO(3)\)}{Full Phase Space: T*SO(3)}}
\label{sec:rigid-full-phase-space}

The Hamiltonian phase space of a finite-dimensional mechanical system with
configuration space \(Q\) is the cotangent bundle \(T^*Q\). Therefore the
unreduced phase space of a rigid body with fixed center of mass is
\[
  T^*\SO(3).
\]
If \(q^i\), \(i=1,2,3\), are any nondegenerate local coordinates on
\(\SO(3)\), the canonical momenta are
\[
  p_i=\frac{\partial L}{\partial \dot q^i}.
\]
Under a change of coordinates \(q^i=q^i(\tilde q)\), the momenta transform as a
one-form:
\[
  \tilde p_j
  =
  \frac{\partial q^i}{\partial \tilde q^j}p_i.
\]
This is the coordinate calculation behind the intrinsic statement
\((q,p)\in T^*Q\).

\begin{remark}[Angular momentum is not a canonical momentum]
Let \(M_a\) be the components of body angular momentum in principal axes:
\[
  M_a=I_a\Omega_a.
\]
The \(M_a\) are excellent coordinates on the fiber of \(T^*\SO(3)\) after a
choice of body frame, but they are not canonical momenta conjugate to Euler
angles. Locally one has a nondegenerate linear relation
\[
  p_i=A_{ia}(q)M_a,
\]
where the matrix \(A(q)\) depends on the chosen angle coordinates. The canonical
symplectic form is still
\[
  \omega=\dd q^i\wedge \dd p_i,
\]
in any cotangent coordinate chart, but when it is expressed in the
noncanonical variables \(M_a\), the Poisson brackets become the rigid-body
Lie-Poisson brackets; the variables \(M_a\) are not canonical coordinates with
\(\Poisson{M_a}{M_b}=0\).
Globally, because \(\SO(3)\) is a Lie group, left or right translation
trivializes the cotangent bundle:
\[
  T^*\SO(3)\simeq \SO(3)\times\mathfrak{so}(3)^*.
\]
This global trivialization is the precise sense in which the angular-momentum
components give global fiber coordinates, even though they are noncanonical.
\end{remark}

\section{\texorpdfstring{Variations On \(\SO(3)\)}{Variations On SO(3)}}
\label{sec:variations-so3}

A variation of \(R(t)\in\SO(3)\) must remain tangent to the group. Write
\[
  \delta R=R\widehat\Sigma,
\]
where \(\Sigma(t)\in\R^3\) vanishes at the time endpoints. The induced variation
of body angular velocity is
\[
  \delta\widehat\Omega
  =
  \delta(R^T\dot R)
  =
  \dot{\widehat\Sigma}+[\widehat\Omega,\widehat\Sigma].
\]
Under the hat isomorphism this becomes
\[
  \delta\Omega=\dot\Sigma+\Omega\times\Sigma.
\]

\begin{derivation}[Euler-Poincare equation for the free rigid body]
This derivation differs from the ordinary Euler-Lagrange calculation because
\(\Omega\) is not an independent velocity coordinate. It is the body-frame
velocity \(R^{-1}\dot R\), so its allowed variation must come from an allowed
variation of \(R\in\SO(3)\). The formula
\(\delta\Omega=\dot\Sigma+\Omega\times\Sigma\) is the constraint that keeps the
variation on the rotation group. Poincare introduced this reduced variational
form in 1901; the modern geometric setting, where it appears as
cotangent-bundle reduction by a Lie group, is the Marsden-Weinstein framework
(1974).
The reduced action is
\[
  S[R]=\int_{t_0}^{t_1}\ell(\Omega)\,\dd t,
  \qquad
  \ell(\Omega)=\frac12\Omega^TI\Omega.
\]
Let
\[
  M=\frac{\partial \ell}{\partial \Omega}=I\Omega.
\]
Then
\[
\begin{aligned}
  \delta S
  &=
  \int_{t_0}^{t_1} M\cdot(\dot\Sigma+\Omega\times\Sigma)\,\dd t\\
  &=
  \left[M\cdot\Sigma\right]_{t_0}^{t_1}
  -
  \int_{t_0}^{t_1}\dot M\cdot\Sigma\,\dd t
  +
  \int_{t_0}^{t_1}M\cdot(\Omega\times\Sigma)\,\dd t.
\end{aligned}
\]
The boundary term vanishes. Since
\[
  M\cdot(\Omega\times\Sigma)=(M\times\Omega)\cdot\Sigma,
\]
stationarity for arbitrary \(\Sigma\) gives
\[
  -\dot M+M\times\Omega=0.
\]
Equivalently,
\[
  \boxed{\dot M+\Omega\times M=0}
\]
with the hat convention \(\widehat\Omega v=\Omega\times v\) used throughout
these notes. In components this is
\[
  I_1\dot\Omega_1=(I_2-I_3)\Omega_2\Omega_3,
\]
and cyclic permutations.
\end{derivation}

The term \(\Omega\times M\) is not an external torque. It is the correction
that appears because \(M\) is expressed in the rotating body frame. In the
space frame the angular momentum \(J=RM\) is constant; in the body frame the
components of the same vector change because the axes themselves rotate.
Euler derived these equations in 1765 in \emph{Theoria motus corporum
solidorum}; the modern reduced variational form is the Euler-Poincare version
of the same calculation.

\section{Euler Equations In Principal Axes}
\label{sec:euler-equations-principal}

With \(M_i=I_i\Omega_i\), Euler's equations are
\[
\begin{aligned}
  I_1\dot\Omega_1 &= (I_2-I_3)\Omega_2\Omega_3,\\
  I_2\dot\Omega_2 &= (I_3-I_1)\Omega_3\Omega_1,\\
  I_3\dot\Omega_3 &= (I_1-I_2)\Omega_1\Omega_2.
\end{aligned}
\]
These equations are autonomous and quadratic. They conserve the kinetic energy
\[
  H=\frac12(I_1\Omega_1^2+I_2\Omega_2^2+I_3\Omega_3^2)
  =
  \frac12\left(\frac{M_1^2}{I_1}+\frac{M_2^2}{I_2}+\frac{M_3^2}{I_3}\right)
\]
and the squared angular momentum
\[
  C=\norm{M}^2=M_1^2+M_2^2+M_3^2.
\]

\begin{proof}[Direct conservation check]
Since \(M=I\Omega\),
\[
  \dot H=\Omega\cdot \dot M.
\]
Using \(\dot M=-\Omega\times M\),
\[
  \dot H=-\Omega\cdot(\Omega\times M)=0.
\]
Similarly,
\[
  \dot C=2M\cdot\dot M=-2M\cdot(\Omega\times M)=0.
\]
\end{proof}

\section{Reconstruction Of The Attitude}
\label{sec:rigid-attitude-reconstruction}

Euler's equations are the reduced equations: they determine the body angular
momentum \(M(t)\) or body angular velocity \(\Omega(t)=I^{-1}M(t)\), but they
do not by themselves tell us the attitude \(R(t)\). Once the reduced solution
is known, the attitude is recovered from the reconstruction equation
\[
  \boxed{
  \dot R=R\widehat\Omega,\qquad R(0)=R_0.
  }
\]
This is a linear equation on the group along the already-computed reduced
motion. In general its solution is a path-ordered exponential,
\[
  R(t)=R_0\,\mathcal P
  \exp\left(\int_0^t\widehat\Omega(s)\,\dd s\right),
\]
where \(\mathcal P\) denotes path ordering, because
\(\widehat\Omega(t_1)\) and \(\widehat\Omega(t_2)\) need not commute.
The solution remains in \(\SO(3)\). Indeed,
\[
  \frac{\dd}{\dd t}(R^TR)
  =
  (R\widehat\Omega)^TR+R^TR\widehat\Omega
  =
  -\widehat\Omega+\widehat\Omega=0,
\]
so \(R^TR\) stays equal to \(I\) if it is \(I\) initially; continuity keeps
\(\det R=1\).

The same equation explains why spatial angular momentum is conserved. Define
\[
  J(t)=R(t)M(t).
\]
Then
\[
\begin{aligned}
  \dot J
  &=
  \dot R M+R\dot M\\
  &=
  R(\Omega\times M+\dot M)\\
  &=
  0
\end{aligned}
\]
by Euler's equation. Thus the reduced equation \(\dot M+\Omega\times M=0\) is
equivalent to the statement that the physical angular momentum vector \(J\) is
fixed in space.

\begin{figure}[htbp]
\centering
\begin{tikzpicture}[scale=1.0, >=Latex, font=\small]
  \draw[->] (-0.2,0) -- (3.25,0) node[below right] {\(M_1\)};
  \draw[->] (0,-0.25) -- (0,3.25) node[above] {\(M_2\)};
  \draw[blue!70!black, thick]
    (1.35,1.55) ellipse (1.20 and 0.58);
  \node[blue!70!black, align=center, fill=white, inner sep=2pt] at (1.05,2.70)
    {Euler flow\\on reduced curve};
  \foreach \x/\y in {0.55/1.55,1.35/2.13,2.15/1.55}
  {
    \fill[blue!70!black] (\x,\y) circle (1.1pt);
    \draw[red!75!black] (\x,\y) circle (0.26);
    \draw[red!75!black, ->] (\x,\y) -- ++(0.26,0.0);
  }
  \node[red!75!black, align=center, fill=white, inner sep=2pt] at (5.25,2.98)
    {attitude fiber\\over a reduced state};
  \draw[red!75!black, ->, thick] (2.55,1.82) .. controls (3.0,2.35) .. (4.0,2.10);
  \draw[red!75!black, ->, thick] (2.45,1.28) .. controls (3.0,0.70) .. (4.0,1.05);
  \begin{scope}[xshift=4.1cm]
    \draw[thick] (0,0.55) rectangle (2.10,2.55);
    \draw[->, purple!75!black, thick]
      (0.18,0.92) .. controls (0.55,1.55) and (1.05,1.72) .. (1.42,2.25);
    \node[purple!75!black, align=center, fill=white, inner sep=2pt] at (2.36,1.05)
      {reconstruct \(R(t)\)\\from \(\dot R=R\widehat\Omega\)};
  \end{scope}
\end{tikzpicture}
\caption{Reduction and reconstruction for the free rigid body. Euler's
equations evolve the body angular momentum on a reduced curve. The attitude is
then reconstructed by integrating \(\dot R=R\widehat\Omega\), with a circle of
possible orientations over a typical reduced state.}
\label{fig:rigid-reduction-reconstruction}
\end{figure}

For numerical work this distinction is important. Integrating \(M\) or
\(\Omega\) checks the reduced invariants \(H\) and \(C=\norm M^2\). Integrating
\(R\) checks the group constraint \(R^TR=I\), \(\det R=1\), and the conservation
of \(J=RM\). The associated Python demo uses an exponential update
\[
  R_{n+1}
  =
  R_n\exp\left(\Delta t\,\widehat\Omega_{n+1/2}\right)
\]
for the reconstruction step, so the attitude stays on \(\SO(3)\) up to roundoff
while the reduced integrator controls the evolution of \(\Omega\).

\section{Geometry Of The Reduced Motion}
\label{sec:rigid-reduced-geometry}

The reduced motion in \(M\)-space lies on the intersection of:
\[
  C(M)=\norm{M}^2=\text{constant}
\]
and
\[
  H(M)=\frac12M^TI^{-1}M=\text{constant}.
\]
Thus it is the intersection of an angular momentum sphere and an energy
ellipsoid.

\begin{figure}[htbp]
\centering
\begin{tikzpicture}[
    scale=1.05,
    >=Latex,
    x={(1cm,0cm)},
    y={(0.42cm,0.27cm)},
    z={(0cm,1cm)},
    line join=round
  ]
  \draw[->, gray!65] (0,0,0) -- (2.75,0,0) node[right] {\(M_1\)};
  \draw[->, gray!65] (0,0,0) -- (0,2.75,0) node[above right] {\(M_2\)};
  \draw[->, gray!65] (0,0,0) -- (0,0,3.15) node[above] {\(M_3\)};

  \draw[red!65!black, thick, domain=0:360, samples=120, variable=\t]
    plot ({2*cos(\t)}, {2*sin(\t)}, {0});
  \draw[red!65!black, densely dashed, domain=0:360, samples=120, variable=\t]
    plot ({2*cos(\t)}, {0}, {2*sin(\t)});
  \draw[red!65!black, loosely dashed, domain=0:360, samples=120, variable=\t]
    plot ({0}, {2*cos(\t)}, {2*sin(\t)});

  \draw[blue!70!black, thick, domain=0:360, samples=140, variable=\t]
    plot ({sqrt(5/2)*cos(\t)}, {sqrt(5)*sin(\t)}, {0});
  \draw[blue!70!black, thin, domain=0:360, samples=140, variable=\t]
    plot ({sqrt(5/2)*cos(\t)}, {0}, {sqrt(15/2)*sin(\t)});
  \draw[blue!70!black, thin, domain=0:360, samples=140, variable=\t]
    plot ({0}, {sqrt(5)*cos(\t)}, {sqrt(15/2)*sin(\t)});

  \draw[purple!80!black, very thick, domain=0:360, samples=180, variable=\t]
    plot ({sqrt(1 + 0.75*sin(\t)*sin(\t))}, {1.732*cos(\t)}, {1.5*sin(\t)});
  \draw[purple!80!black, very thick, dashed, domain=0:360, samples=180, variable=\t]
    plot ({-sqrt(1 + 0.75*sin(\t)*sin(\t))}, {1.732*cos(\t)}, {1.5*sin(\t)});

  \fill[purple!80!black] ({sqrt(1 + 0.75*0*0)}, {1.732}, {0}) circle (1.3pt);
  \node[red!65!black, align=center, fill=white, inner sep=2pt] at (-2.65,-1.25,-0.75)
    {angular momentum\\sphere \(C=\norm M^2\)};
  \node[blue!70!black, align=center, fill=white, inner sep=2pt] at (2.75,1.75,3.28)
    {energy ellipsoid\\\(2H=M^TI^{-1}M\)};
  \node[purple!80!black, fill=white, inner sep=2pt] at (4.65,-2.75,-0.55)
    {intersection trajectories};
  \draw[purple!80!black, ->] (3.55,-2.48,-0.45) -- (1.15,-0.25,0);
  \node[purple!80!black, align=center, fill=white, inner sep=1.5pt] at (-3.95,2.25,0.40)
    {back branch\\(dashed)};
  \draw[purple!80!black, ->] (-3.30,1.95,0.35) -- (-1.05,0.95,0.55);
\end{tikzpicture}
\caption{Rigid-body reduction in angular-momentum space. Conservation of
\(\norm M^2\) confines the motion to a sphere, conservation of \(H\) confines it
to an ellipsoid, and the reduced trajectory is their intersection. The solid
and dashed purple curves are the two projected branches of the exact
three-dimensional intersection for the sphere \(\norm M^2=4\) and the
representative ellipsoid with semi-axes
\(\sqrt{5/2},\sqrt5,\sqrt{15/2}\).}
\label{fig:rigid-body-reduced-geometry}
\end{figure}

\section{Poinsot Construction}
\label{sec:poinsot-construction}

The same invariants have a classical geometric interpretation in angular
velocity space. Poinsot introduced this rolling picture in 1834. His vocabulary
survives: the body-frame trace is the polhode, and the spatial trace is the
herpolhode. Define the inertia ellipsoid
\[
  \mathcal E_H=\{\Omega:\Omega^TI\Omega=2H\}.
\]
At a point \(\Omega\) on this ellipsoid, the normal vector is proportional to
\[
  \nabla_\Omega(\Omega^TI\Omega)=2I\Omega=2M.
\]
In spatial coordinates the corresponding angular velocity is
\[
  \omega_{\mathrm{space}}=R\Omega,
\]
and the normal is the fixed spatial angular momentum
\[
  J=RM.
\]
Because
\[
  J\cdot\omega_{\mathrm{space}}
  =
  (RM)\cdot(R\Omega)
  =
  M\cdot\Omega
  =
  2H,
\]
the tip of \(\omega_{\mathrm{space}}\) lies in the fixed plane
\[
  \boxed{\Pi_J=\{\omega:J\cdot\omega=2H\}.}
\]
This plane is perpendicular to \(J\) and has distance \(2H/\norm J\) from the
origin. It is tangent to the moving inertia ellipsoid at the instantaneous
angular velocity. Indeed, the tangent plane to
\(\Omega^TI\Omega=2H\) has normal \(M=I\Omega\); after rotation to space the
normal is \(J=RM\), the normal of \(\Pi_J\). The point of contact is
\(\omega_{\mathrm{space}}=R\Omega\), and it lies in \(\Pi_J\) by the displayed
identity \(J\cdot\omega_{\mathrm{space}}=2H\).

The no-slip statement is also kinematic. The instantaneous axis of the rigid
motion is \(\omega_{\mathrm{space}}\), so the point of the moving ellipsoid
whose spatial position is the contact point has zero instantaneous velocity
relative to the fixed plane. Thus the free rigid body can be pictured as the
inertia ellipsoid rolling without slipping on the fixed invariable plane.

\begin{figure}[htbp]
\centering
\begin{tikzpicture}[scale=1.0, >=Latex]
  \begin{scope}[xshift=-0.4cm]
    \draw[fill=blue!4, draw=blue!70!black, thick] (0,0) ellipse (1.65 and 0.85);
    \draw[blue!70!black, dashed] (-1.65,0) arc (180:360:1.65 and 0.28);
    \draw[blue!70!black] (1.65,0) arc (0:180:1.65 and 0.28);
    \draw[purple!75!black, thick]
      (-0.9,0.15) .. controls (-0.35,0.85) and (0.7,0.55) .. (1.05,-0.05);
  \fill[purple!75!black] (0.18,0.55) circle (1.5pt);
  \node[purple!75!black, fill=white, inner sep=1pt] at (-0.08,0.82) {\(\Omega\)};
    \draw[->, red!75!black, thick] (0.18,0.55) -- (1.05,1.32);
    \node[red!75!black, fill=white, inner sep=1pt] at (1.25,1.42)
      {\(M=I\Omega\)};
    \node[blue!70!black] at (0,-1.25) {inertia ellipsoid \(\mathcal E_H\)};
    \node[purple!75!black] at (-1.25,0.95) {polhode};
  \end{scope}
  \begin{scope}[xshift=4.6cm, yshift=0.1cm]
    \draw[fill=gray!12, draw=gray!75, thick]
      (-1.7,-0.55) -- (1.8,-0.15) -- (1.2,0.95) -- (-2.1,0.55) -- cycle;
    \draw[->, red!75!black, thick] (0,0.15) -- (0,2.0)
      node[above] {fixed \(J\)};
    \draw[purple!75!black, thick]
      (-1.0,0.05) .. controls (-0.25,0.55) and (0.9,0.42) .. (1.25,0.18);
    \fill[purple!75!black] (0.05,0.43) circle (1.5pt)
      node[below right] {\(\omega_{\mathrm{space}}\)};
    \node[gray!70!black] at (0,-0.95) {invariable plane \(\Pi_J\)};
    \node[purple!75!black] at (1.35,0.75) {herpolhode};
  \end{scope}
\end{tikzpicture}
\caption{Poinsot construction. In body angular-velocity space, the motion lies
on the inertia ellipsoid and traces the polhode. In space, the tip of angular
velocity lies in the fixed invariable plane perpendicular to the conserved
spatial angular momentum \(J\); its trace is the herpolhode.}
\label{fig:poinsot-construction}
\end{figure}

Poinsot's construction is another geometric view of the same conservation of
\(H\) and \(J\), drawn in angular-velocity space. Its value is conceptual: it
makes clear why the reduced trajectory is a curve, why the invariable plane is
fixed in space, and why body-frame quantities can move while the spatial
angular momentum does not.

\section{Stability Of Principal-Axis Rotations}
\label{sec:principal-axis-stability}

Assume
\[
  I_1<I_2<I_3.
\]
A steady rotation about a principal axis is a solution. For example,
\[
  \Omega(t)=(\Omega_0,0,0)
\]
is rotation about the first axis.

\begin{derivation}[Linear stability]
Perturb rotation about the first axis:
\[
  \Omega_1=\Omega_0+\alpha_1,\qquad
  \Omega_2=\alpha_2,\qquad
  \Omega_3=\alpha_3,
\]
and keep only linear terms. Then
\[
  \dot\alpha_1=0,
\]
while
\[
  I_2\dot\alpha_2=(I_3-I_1)\Omega_0\alpha_3,
  \qquad
  I_3\dot\alpha_3=(I_1-I_2)\Omega_0\alpha_2.
\]
Differentiate the first of these:
\[
  \ddot\alpha_2
  =
  \frac{(I_3-I_1)(I_1-I_2)}{I_2I_3}\Omega_0^2\alpha_2.
\]
Since \(I_3-I_1>0\) and \(I_1-I_2<0\), the coefficient is negative, so
\(\alpha_2\) oscillates. Rotation about the first axis is linearly stable.

For rotation about the second axis, the analogous coefficient is
\[
  \frac{(I_2-I_3)(I_1-I_2)}{I_1I_3}\Omega_0^2>0,
\]
because both factors in the numerator are negative. The perturbation therefore
has exponential solutions. Thus the intermediate axis is unstable. Rotation
about the third axis is stable.
\end{derivation}

The geometric reason is that the energy restricted to a fixed angular momentum
sphere has extrema at the smallest and largest inertia axes but a saddle at the
intermediate axis.
This instability is the tennis-racket theorem; it was filmed by cosmonaut
Vladimir Dzhanibekov on Salyut-7 in 1985. The mathematics is already in
Euler's 1765 rigid-body equations.

\begin{derivation}[Energy-Casimir stability test]
The same conclusion can be read without solving the linearized equations. Work
on the angular-momentum sphere
\[
  C(M)=M_1^2+M_2^2+M_3^2=c^2.
\]
Near rotation about the first principal axis, write
\[
  M=(M_1,\eta_2,\eta_3),
  \qquad
  M_1=\sqrt{c^2-\eta_2^2-\eta_3^2}.
\]
Expanding the Hamiltonian to second order gives
\[
\begin{aligned}
  H-H_0
  &=
  \frac12\left(\frac{1}{I_2}-\frac{1}{I_1}\right)\eta_2^2
  +
  \frac12\left(\frac{1}{I_3}-\frac{1}{I_1}\right)\eta_3^2
  +
  O(\eta^4).
\end{aligned}
\]
If \(I_1<I_2<I_3\), both quadratic coefficients are negative. The first-axis
rotation is therefore a strict local maximum of \(H\) on the sphere \(C=c^2\).
The third-axis rotation is similarly a strict local minimum. Either extremum
prevents nearby trajectories with the same \(C\) and nearly the same \(H\) from
leaving a small neighborhood.

For the intermediate axis,
\[
  M=(\eta_1,M_2,\eta_3),
\]
and
\[
  H-H_0
  =
  \frac12\left(\frac{1}{I_1}-\frac{1}{I_2}\right)\eta_1^2
  +
  \frac12\left(\frac{1}{I_3}-\frac{1}{I_2}\right)\eta_3^2
  +
  O(\eta^4).
\]
The two coefficients have opposite signs, so the point is a saddle of the
energy restricted to the angular-momentum sphere. This is the energy-Casimir
form of the intermediate-axis instability.
\end{derivation}

\section{Lie-Poisson Structure}
\label{sec:lie-poisson-rigid-body}

The reduced rigid body is Hamiltonian on \(\mathfrak{so}(3)^*\cong\R^3\), but
not with canonical coordinates. The Lie-Poisson bracket is
\[
  \Poisson{F}{G}(M)
  =
  -M\cdot\left(\nabla F(M)\times\nabla G(M)\right).
\]
In coordinate functions this says
\[
  \Poisson{M_i}{M_j}=-\epsilon_{ijk}M_k.
\]
This is the minus-Lie-Poisson convention. The minus sign matches the
body-frame variables used above and gives the boxed Euler equation
\(\dot M=-\Omega\times M\). Had we used spatial angular momentum components as
reduced variables, the sign would be opposite. This is another place where the
body/space distinction is not cosmetic: it is the \(\pm\) Lie-Poisson
dichotomy coming from left versus right trivialization of \(T^*\SO(3)\). The
physical motion is the same, but the reduced bracket changes sign.

For
\[
  H(M)=\frac12M\cdot I^{-1}M,
  \qquad
  \nabla H=I^{-1}M=\Omega,
\]
the equation for any observable \(F\) is
\[
  \dot F=\Poisson{F}{H}.
\]
Taking \(F(M)=M_i\) yields
\[
  \dot M=-\Omega\times M.
\]

\begin{definition}[Casimir]
A Casimir is a function \(C\) satisfying
\[
  \Poisson{C}{F}=0
\]
for every observable \(F\). For the rigid body bracket,
\[
  C(M)=\norm{M}^2
\]
is a Casimir. Its level sets are the coadjoint orbits, which here are spheres.
\end{definition}

Casimirs foliate the reduced phase space. For the rigid body,
\(\norm M^2\) labels each angular-momentum sphere, and the Hamiltonian flow
lives on these spheres. Each sphere carries a genuine two-dimensional
symplectic structure, and the rigid body restricted to one such sphere has one
degree of freedom.
Here a coadjoint orbit means the orbit of a covector
\(M\in\mathfrak{so}(3)^*\) under the coadjoint action of \(\SO(3)\). After
identifying \(\mathfrak{so}(3)^*\) with \(\R^3\), that action is ordinary
rotation, so the nonzero orbits are spheres \(\norm M=\mathrm{constant}\).
This is the rigid-body instance of Marsden-Weinstein symplectic reduction
(1974): the canonical phase space \(T^*\SO(3)\) descends to coadjoint orbits
with their Kirillov-Kostant-Souriau symplectic forms.
This is the reduced reason that the intersection curves in
Figure~\ref{fig:rigid-body-reduced-geometry} are the phase curves of the
system.

\section{Liouville Integrability And Invariant Tori}
\label{sec:liouville-tori}

The rigid-body calculations above show a pattern: conservation laws confine
motion to lower-dimensional sets, and the reduced flow on those sets can be
simple even when the original coordinates look complicated. Liouville
integrability is the precise Hamiltonian version of this pattern. It is not
just ``having many constants of motion.'' The constants must be independent,
they must Poisson commute, and their common level set must be regular enough to
carry the commuting Hamiltonian flows.

The three hypotheses have separate meanings. Independence says the level set
has the expected dimension. Involution says the Hamiltonian flows generated by
the integrals commute. Compactness says the flow coordinates eventually repeat,
turning straight-line motion on \(\R^n\) into straight-line motion on a torus.
Connectedness says one is discussing a single component of the common level
set; otherwise the theorem applies component by component. Completeness of the
commuting flows is the dynamical hypothesis that lets the \(\R^n\)-action be
followed for all flow times. Dropping any one of these assumptions changes the
conclusion.

This is why the theorem should not be read backwards. Seeing a bounded-looking
orbit or many conserved-looking quantities in a simulation does not by itself
prove integrability. One must know that the integrals are globally defined on
the phase-space region under study, independent on the level set, and in
involution. The free rigid body passes these tests; a slightly forced or
gravity-coupled top often fails them.

\begin{definition}[Liouville integrability]
A Hamiltonian system on a \(2n\)-dimensional symplectic manifold is Liouville
integrable if there are \(n\) independent functions
\[
  F_1=H,\ F_2,\ldots,F_n
\]
such that
\[
  \Poisson{F_i}{F_j}=0
\]
for all \(i,j\), and the associated Hamiltonian flows are complete on the
regular level sets under discussion. The functions are said to be in
involution.
\end{definition}

Liouville proved the analytic integration-by-quadratures part in 1855: once
enough commuting integrals are present, the remaining motion can be obtained by
quadratures. The modern compact-level-set formulation, due to Arnold (1963), is
now called the Liouville-Arnold theorem; it identifies regular compact level
sets as tori carrying linear flow.
Jacobi's 1839 integration of geodesics on an ellipsoid is a classical early
example beyond Kepler and the rigid body: the separated coordinates are not a
coordinate trick but a sign of hidden commuting integrals.

\begin{example}[Basic integrable systems]
Any autonomous Hamiltonian system with one degree of freedom is integrable: the
Hamiltonian itself supplies the one conserved quantity. The fixed-COM rigid body
has three degrees of freedom. The three spatial functions
\[
  H,\qquad J_3,\qquad J^2
\]
are independent and mutually Poisson commuting on regular sets, even though the
individual components \(J_1,J_2,J_3\) do not commute with one another. This is
the rigid-body example used throughout this chapter.
\end{example}

\begin{assumption}[Regular compact level set]
Fix constants \(f=(f_1,\ldots,f_n)\) and define
\[
  M_f=\{x\in P: F_i(x)=f_i,\ i=1,\ldots,n\},
\]
where \(P\) is the \(2n\)-dimensional phase space. We assume that the
differentials \(\dd F_i\) are independent on \(M_f\) and that \(M_f\) is
compact and connected.
\end{assumption}

The regularity assumption says that \(M_f\) is a smooth \(n\)-dimensional
submanifold. Its tangent vectors \(v\in T_xM_f\) are those satisfying
\[
  \dd F_i(v)=0
  \qquad\text{for all }i.
\]
Let \(X_{F_i}\) be the Hamiltonian vector field of \(F_i\), defined by
\[
  \iota_{X_{F_i}}\omega=\dd F_i.
\]
Then
\[
  \dd F_j(X_{F_i})
  =
  \Poisson{F_j}{F_i}
  =
  0.
\]
Thus every \(X_{F_i}\) is tangent to \(M_f\). The commutator of two Hamiltonian
vector fields is
\[
  [X_{F_i},X_{F_j}]
  =
  -X_{\Poisson{F_i}{F_j}},
\]
so involution implies
\[
  [X_{F_i},X_{F_j}]=0.
\]
We therefore get \(n\) independent commuting vector fields on \(M_f\).

\begin{derivation}[Why the regular compact level set is a torus]
The proof has a simple geometric shape. The Hamiltonian vector fields
\(X_{F_i}\) give \(n\) independent directions tangent to \(M_f\), and because
the integrals Poisson commute, flowing in one direction and then another gives
the same point as flowing in the opposite order. Thus \(M_f\) is locally
coordinatized by ``flow times.'' Compactness is what turns these flow-time
coordinates from \(\R^n\) into periodic coordinates.

Commuting vector fields have commuting flows. Starting at a point \(x_0\in
M_f\), flow for times \(s=(s^1,\ldots,s^n)\) along the \(n\) vector fields:
\[
  \Phi_s(x_0)
  =
  \exp(s^1X_{F_1})\cdots\exp(s^nX_{F_n})x_0.
\]
Because the flows commute, this defines an action of \(\R^n\) on \(M_f\).
Completeness is used here: without complete flows, this action would
only be defined for small flow times, and the period-lattice quotient below
would have no global meaning. Since \(M_f\) is compact and the vector fields
are smooth, the flows on \(M_f\) are complete in the cases considered here.
This is the standard escape lemma for ODEs: a maximal solution of a smooth
vector field that remained defined only on a finite time interval would have
compact closure, the vector field would be bounded there, and the local
existence theorem would extend the solution past the endpoint.
Independence of the vector fields makes the action locally free. More
concretely, the constant-rank theorem applied to the map
\(s\mapsto\Phi_s(x_0)\) says that its derivative at \(s=0\) is an isomorphism
onto \(T_{x_0}M_f\); hence the flow times \(s^i\) are local coordinates. The
only possible global identifications are translations
\[
  s\sim s+\lambda,\qquad \lambda\in\Lambda,
\]
where \(\Lambda\subset\R^n\) is the period lattice of the orbit. Compactness
forces \(\Lambda\) to have rank \(n\). If \(\Lambda\) had rank \(k<n\), then
the structure theorem for discrete subgroups of \(\R^n\) gives a linear basis
in which \(\Lambda\) is generated by \(k\) independent period vectors, so
\[
  \R^n/\Lambda\simeq \mathbb T^k\times\R^{n-k},
\]
which is noncompact. The action is locally free, so each orbit is open in the
\(n\)-manifold \(M_f\). Since \(M_f\) is connected, that open orbit is the
whole level set. Hence
\[
  M_f\simeq \R^n/\Lambda\simeq \mathbb T^n.
\]
These are the invariant Liouville tori.
\end{derivation}

Enough commuting flows on a compact regular level set force the level set to be
a torus. The angle variables introduced below are coordinates
along these flows, and the action variables are the quantities that label which
torus one is on.

\begin{figure}[htbp]
\centering
\begin{tikzpicture}[scale=1.0, >=Latex, font=\small]
  \draw[thick] (0,0) rectangle (4.2,3.0);
  \draw[->, blue!70!black, thick] (0.45,0.55) -- (3.75,0.55)
    node[midway, below] {\(\gamma_1\)};
  \draw[->, red!75!black, thick] (0.55,0.45) -- (0.55,2.55)
    node[midway, left] {\(\gamma_2\)};
  \draw[purple!75!black, thick, ->]
    (0.25,0.25) -- (1.1,0.9) -- (1.95,1.55) -- (2.8,2.2) -- (3.65,2.85);
  \draw[purple!75!black, thick, ->]
    (3.65,2.85) -- (4.15,3.23);
  \draw[purple!75!black, thick, ->]
    (-0.15,-0.23) -- (0.25,0.25);
  \draw[gray!70, ->] (4.2,1.5) -- (4.8,1.5)
    node[right, fill=white, inner sep=1.5pt] {identify};
  \draw[gray!70, ->] (2.1,3.0) -- (2.1,3.55)
    node[above, fill=white, inner sep=1.5pt] {identify};
  \node[align=center, fill=white, inner sep=1.5pt] at (2.1,-0.72)
    {flat model of a two-dimensional invariant torus\\
     \(\mathbb R^2/\Lambda\) with angle coordinates};
  \node[purple!75!black, align=center, fill=white, inner sep=2pt] at (5.05,3.18)
    {Hamiltonian flow\\has constant slope};
  \draw[purple!75!black, ->] (4.70,2.98) -- (3.20,2.48);
\end{tikzpicture}
\caption{A Liouville torus is locally a space of commuting flow times with
periodic identifications. Action variables label the torus; angle variables
are coordinates along the fundamental cycles.}
\label{fig:liouville-flat-torus}
\end{figure}

The symplectic form vanishes when restricted to an invariant torus. Indeed, on
the tangent basis \(X_{F_i}\),
\[
  \omega(X_{F_i},X_{F_j})
  =
  \dd F_i(X_{F_j})
  =
  \Poisson{F_i}{F_j}
  =
  0.
\]
Since these vectors span \(T_xM_f\), the torus is Lagrangian:
\[
  \omega|_{M_f}=0.
\]

\begin{derivation}[Coordinate proof that the invariant torus is Lagrangian]
A useful coordinate check is the following. Use canonical
coordinates \((q_i,p_i)\), and suppose the level set is written locally as
\[
  F_i(p,q)=f_i,\qquad i=1,\ldots,n.
\]
If the equations can be solved locally for
\[
  p_i=p_i(q,F),
\]
then
\[
\begin{aligned}
  \omega
  &=
  \dd q_i\wedge\dd p_i\\
  &=
  \dd q_i\wedge
  \left(
    \frac{\partial p_i}{\partial q_j}\bigg|_F\dd q_j
    +
    \frac{\partial p_i}{\partial F_j}\bigg|_q\dd F_j
  \right).
\end{aligned}
\]
The key claim is
\[
  \frac{\partial p_i}{\partial q_j}\bigg|_F
  =
  \frac{\partial p_j}{\partial q_i}\bigg|_F.
\]
Then the \(\dd q_i\wedge\dd q_j\) contribution vanishes by antisymmetry, and
\[
  \omega=
  \frac{\partial p_i}{\partial F_j}\bigg|_q\dd q_i\wedge\dd F_j.
\]
On the invariant torus \(F_j=f_j\), \(\dd F_j=0\), so
\[
  \omega|_{M_f}=0.
\]
This coordinate proof is local on a chart where the projection of the torus to
the \(q\)-coordinates has full rank. At caustics or turning sets that
projection fails, the same Lagrangian conclusion remains true by the
coordinate-free argument above: at a regular point the independent Hamiltonian
fields \(X_{F_i}\) still span \(T_xM_f\), and their pairwise symplectic
products are still \(\Poisson{F_i}{F_j}=0\). What fails at a caustic is only
the choice of \(q\) as local coordinates on the torus. The displayed function
\(p(q,F)\) must then be replaced by another local branch or another canonical
chart.

To prove the claim, differentiate the identities \(F_k(p(F,q),q)=F_k\) at
fixed \(F\):
\[
  0=
  \frac{\partial F_k}{\partial p_j}\bigg|_q
  \frac{\partial p_j}{\partial q_i}\bigg|_F
  +
  \frac{\partial F_k}{\partial q_i}\bigg|_p.
\]
Let
\[
  A_{kj}=\frac{\partial F_k}{\partial p_j}\bigg|_q,\qquad
  M_{ji}=\frac{\partial p_j}{\partial q_i}\bigg|_F,\qquad
  B_{ki}=\frac{\partial F_k}{\partial q_i}\bigg|_p.
\]
Then
\[
  AM+B=0,\qquad M=-A^{-1}B.
\]
The involution condition gives
\[
  0=\Poisson{F_i}{F_j}
  =
  B_{ik}A_{jk}-A_{ik}B_{jk},
\]
or, equivalently,
\[
  AB^T-BA^T=0.
\]
Hence
\[
  A^{-1}B=B^T(A^T)^{-1}=(A^{-1}B)^T,
\]
so \(M=M^T\), proving the symmetry claim.
\end{derivation}

\section{Action Variables And Angle Variables}
\label{sec:action-angle-variables}

Choose a basis of one-cycles \(\gamma_1,\ldots,\gamma_n\) on an invariant torus
\(M_f\). In canonical coordinates the action variables are
\[
  I_j=\frac{1}{2\pi}\oint_{\gamma_j} p_i\,\dd q^i.
\]
Actions are therefore not arbitrary constants of integration. They are
phase-space areas measured around the independent cycles of the torus. This is
why they are insensitive to small deformations of the cycle on a fixed torus.
They are not unique, however: changing the integer cycle basis of
\(H_1(\mathbb T^n;\mathbb Z)\) transforms the action vector by an integral
unimodular matrix, \(I\mapsto AI\) with \(A\in\operatorname{GL}(n,\mathbb Z)\) (or
\(\operatorname{SL}(n,\mathbb Z)\) after fixing orientation). Thus a local choice of action
variables is part of the action-angle chart, not a canonical object independent
of cycle basis.
They are invariant under continuous deformations of \(\gamma_j\) on the torus.
If \(\gamma_j\) and \(\gamma_j'\) bound a two-chain \(D\subset M_f\), then by
Stokes' theorem
\[
  \oint_{\gamma_j}p_i\,\dd q^i
  -
  \oint_{\gamma_j'}p_i\,\dd q^i
  =
  \int_D \dd(p_i\,\dd q^i)
  =
  \int_D \omega
  =
  0.
\]

This is also why action variables entered early quantum theory: the
Bohr-Sommerfeld quantization rules, especially Sommerfeld's 1916 old-quantum
conditions, imposed \(\oint p\,\dd q=nh\) on precisely these periods. The modern
classical use is not quantization, but the same periods are the coordinates that
survive canonical perturbation theory; in one degree of freedom, the relation
between action and period is also the classical seed of WKB quantization.
Thus \(I_j\) depends on the torus, not on the representative loop. Under
regularity assumptions, the \(I_j\) can replace the original integrals \(F_j\)
as labels for nearby tori.

To find conjugate angles, construct a local generating function
\[
  S(I,q)
  =
  \int^q p_i(I,q')\,\dd q'^i,
\]
where the integration is performed on a local branch. The compatibility
condition for this definition is
\[
  \frac{\partial p_i}{\partial q^j}\bigg|_I
  =
  \frac{\partial p_j}{\partial q^i}\bigg|_I,
\]
which is another form of \(\omega|_{M_I}=0\). Define
\[
  \theta^j=\frac{\partial S}{\partial I_j}\bigg|_q.
\]
This generating function proves that the new coordinates are canonical. Since
\[
  \dd S=p_i\,\dd q^i+\theta^j\,\dd I_j,
\]
taking an exterior derivative gives
\[
  0=\dd p_i\wedge\dd q^i+\dd\theta^j\wedge\dd I_j.
\]
Thus
\[
  \dd q^i\wedge\dd p_i=\dd\theta^j\wedge\dd I_j,
\]
so the symplectic form takes canonical action-angle form.
When one goes around the cycle \(\gamma_k\), the multivalued function \(S\)
jumps by
\[
  \Delta_{\gamma_k}S=\oint_{\gamma_k}p_i\,\dd q^i=2\pi I_k.
\]
Therefore
\[
  \Delta_{\gamma_k}\theta^j
  =
  \frac{\partial}{\partial I_j}(2\pi I_k)
  =
  2\pi\delta^j{}_k.
\]
The \(\theta^j\) are angles on the torus:
\[
  \theta^j\sim \theta^j+2\pi.
\]
In these action-angle coordinates the Hamiltonian is a function of the actions
alone,
\[
  H=H(I),
\]
and Hamilton's equations reduce to
\[
  \dot I_j=0,\qquad
  \dot\theta^j=\omega^j(I)=\frac{\partial H}{\partial I_j}.
\]
Thus motion on each invariant torus is linear. If the frequencies are rationally
related, the orbit closes; otherwise the orbit is dense on the torus.

The word ``linear'' refers to the angle coordinates, not to the original
configuration variables. A dense quasiperiodic curve on a torus can look quite
complicated after being projected to ordinary coordinates. Action-angle
variables are valuable precisely because they separate the simple invariant
motion on the torus from the sometimes complicated embedding of that torus in
phase space.

\section{Action Integrals In One Degree Of Freedom}
\label{sec:one-dimensional-action-integrals}

The word action now appears in a second, related sense. Hamilton's principle
uses the time integral \(\int L\,\dd t\) as a functional on paths. Action-angle
theory uses
\[
  I_j=\frac{1}{2\pi}\oint_{\gamma_j}p_i\,\dd q^i
\]
as a coordinate labelling a closed orbit or invariant torus. The
common ingredient is the canonical one-form \(p_i\,\dd q^i\), whose periods are
the action variables.

The abstract definition
\[
  I_j=\frac{1}{2\pi}\oint_{\gamma_j}p_i\,\dd q^i
\]
is easiest to understand first in one degree of freedom. Suppose an autonomous
system has Hamiltonian
\[
  H(q,p)=\frac{p^2}{2m}+V(q)
\]
and a bounded orbit between turning points \(q_-(E)\) and \(q_+(E)\). On the
energy curve \(H=E\),
\[
  p(q,E)=\pm\sqrt{2m(E-V(q))}.
\]
The closed phase-space loop has a positive-momentum branch and a
negative-momentum branch, so the action is
\[
  \boxed{
  I(E)=\frac{1}{2\pi}\oint p\,\dd q
  =
  \frac1{\pi}\int_{q_-(E)}^{q_+(E)}
  \sqrt{2m(E-V(q))}\,\dd q.
  }
\]

\begin{figure}[htbp]
\centering
\begin{tikzpicture}[scale=1.0, >=Latex]
  \draw[->] (-0.4,0) -- (6.6,0) node[right] {\(q\)};
  \draw[->] (0,-2.0) -- (0,2.2) node[above] {\(p\)};
  \draw[densely dotted] (1.0,-1.7) -- (1.0,1.8);
  \draw[densely dotted] (5.3,-1.7) -- (5.3,1.8);
  \node[below left=6pt] at (1.0,0) {\(q_-\)};
  \node[below right=6pt] at (5.3,0) {\(q_+\)};
  \draw[blue!70!black, thick, ->]
    (1.0,0) .. controls (2.0,1.55) and (4.2,1.55) .. (5.3,0);
  \draw[blue!70!black, thick, ->]
    (5.3,0) .. controls (4.2,-1.55) and (2.0,-1.55) .. (1.0,0);
  \node[blue!70!black, fill=white, inner sep=1pt] at (3.2,1.95) {\(+p(q,E)\)};
  \node[blue!70!black, fill=white, inner sep=1pt] at (3.2,-2.12) {\(-p(q,E)\)};
  \draw[blue!70!black, ->] (3.2,1.78) -- (3.2,1.40);
  \draw[blue!70!black, ->] (3.2,-1.78) -- (3.2,-1.40);
  \node[align=center, fill=white, inner sep=1pt] at (3.2,0.55)
    {enclosed area\\\(\displaystyle \oint p\,\dd q=2\pi I\)};
\end{tikzpicture}
\caption{For a one-dimensional bound orbit the action is the signed
phase-space area enclosed by the energy curve divided by \(2\pi\). The turning
points contribute no endpoint term when differentiating with respect to
energy because \(p(q_\pm,E)=0\).}
\label{fig:action-loop}
\end{figure}

The period follows by differentiating \(I(E)\). The moving endpoint terms
vanish because \(p(q_\pm(E),E)=0\):
\[
  \frac{\dd I}{\dd E}
  =
  \frac1{\pi}\int_{q_-}^{q_+}
  \frac{m\,\dd q}{\sqrt{2m(E-V(q))}}
  =
  \frac1{2\pi}T(E).
\]
Therefore
\[
  \boxed{
  \omega(E)=\frac{\dd H}{\dd I}
  =
  \left(\frac{\dd I}{\dd E}\right)^{-1},
  \qquad
  T(E)=2\pi\frac{\dd I}{\dd E}.
  }
\]
This formula is the one-dimensional seed of action-angle theory. The
same period integral reappears on higher-dimensional invariant tori by
integrating the canonical one-form around a basis of cycles.

\begin{example}[Harmonic oscillator]
For
\[
  H=\frac{p^2}{2m}+\frac12m\omega_0^2q^2,
\]
the phase curve is an ellipse:
\[
  \frac{p^2}{2mE}+\frac{q^2}{2E/(m\omega_0^2)}=1.
\]
The area enclosed by the ellipse is
\[
  \pi\sqrt{2mE}\sqrt{\frac{2E}{m\omega_0^2}}
  =
  \frac{2\pi E}{\omega_0}.
\]
Thus
\[
  I=\frac{1}{2\pi}\operatorname{Area}=\frac{E}{\omega_0},
  \qquad
  H(I)=\omega_0 I.
\]
The conjugate angle advances uniformly:
\[
  \dot\theta=\frac{\partial H}{\partial I}=\omega_0.
\]
\end{example}

\section{Rigid Body Tori}
\label{sec:rigid-body-tori}

The full rigid body with fixed center of mass has three degrees of freedom. On
\(T^*\SO(3)\) the functions
\[
  H,\qquad J_3,\qquad J^2
\]
are mutually Poisson commuting when \(J_3\) is a fixed spatial component of
angular momentum and \(J^2=\norm{J}^2\). For regular values, their common level
set is a three-torus.

The phrase ``regular values'' hides the familiar exceptional motions. Principal
axis rotations and separatrix-level motions are not generic tori: the level set
can pinch, a cycle can collapse, or the period can diverge. Liouville
integrability describes the open dense region of phase space where the
commuting flows are independent and compact; the exceptional sets are precisely
where the qualitative changes in rigid-body motion occur.

The geometry can be seen in two stages. First, the body angular momentum
\(M=R^TJ\) is constrained by
\[
  J^2=M_1^2+M_2^2+M_3^2=\text{constant}
\]
and
\[
  H=\frac12\left(\frac{M_1^2}{I_1}
  +\frac{M_2^2}{I_2}
  +\frac{M_3^2}{I_3}\right)=\text{constant}.
\]
This puts \(M\) on the closed curve obtained by intersecting a sphere with an
ellipsoid. Second, fixing \(J_3\) constrains the orientation \(R\): the third
spatial component of \(J=RM\) is fixed. For a given \(M\) on the curve, this
condition leaves a circle of allowed attitudes. A circle fibered over the
momentum-space circle gives a two-dimensional torus in the reduced picture, and
including the remaining cyclic angle gives the full three-torus.

The reduced free rigid body has one effective degree of freedom on each angular
momentum sphere. The energy level on that sphere gives closed curves except at
separatrices. This is the concrete meaning of integrability here.

\section{Rigid Body Kinematics As A Gauge Prototype}
\label{sec:rigid-to-deforming}

The rigid body has one orientation \(R(t)\). A Cosserat body assigns an
orientation to each material point:
\[
  R(t)
  \quad\longrightarrow\quad
  R(X,t).
\]
The rigid-body angular velocity
\[
  R^T\dot R
\]
becomes a time connection, while the spatial derivatives
\[
  R^T\partial_iR
\]
become material connection components. The same operation that turns an attitude
matrix into angular velocity therefore reappears, point by point, as the
connection language of deformable bodies.

\section{What To Take Forward}
\label{sec:rigid-take-forward}

\begin{itemize}
\item A rigid body lives on \(\SO(3)\), not in a global set of Euler angles.
Body and spatial angular momentum describe the same vector in different frames,
and the sign of the reduced Lie-Poisson bracket depends on that choice.
\item The free rigid body is integrable because energy and angular momentum
confine the reduced motion to curves on angular-momentum spheres; after
reconstruction these curves lift to invariant tori.
\item Liouville tori are produced by commuting Hamiltonian flows, with action
variables defined by period integrals of the canonical one-form. The operation
\(R^T\dd R\) used here will reappear as the connection language of deformable
bodies and Cosserat continua.
\end{itemize}

\section{Chapter-End Laboratories And Exercises}
\label{sec:rigid-integrability-chapter-end-labs}

These calculations connect action variables to concrete integrable systems and
show how the rigid body fits into the same torus-based language.

\subsection[Oscillator Actions In Polar Coordinates]{Action Variables For The
Isotropic Oscillator In Polar Coordinates}
\label{sec:detailed-oscillator-actions}

A canonical test case for action-angle coordinates is the two-dimensional
isotropic oscillator
\[
  H=
  \frac{p_r^2}{2m}
  +
  \frac{p_\theta^2}{2mr^2}
  +
  \frac12m\omega^2r^2.
\]
The angular momentum
\[
  L=p_\theta
\]
is conserved. The angular action is immediate:
\[
  I_\theta=\frac1{2\pi}\oint p_\theta\,\dd\theta=L.
\]
The radial action is
\[
  I_r
  =
  \frac1{\pi}
  \int_{r_-}^{r_+}
  \sqrt{
    2mE-\frac{L^2}{r^2}-m^2\omega^2r^2
  }\,\dd r.
\]
Rather than evaluate this integral by brute force, use a structural argument.
For the Cartesian oscillator,
\[
  H=\omega(I_x+I_y).
\]
The angular momentum generates rotations in the plane. For nonnegative \(L\),
the circular orbit has all its action in rotation and no radial oscillation.
On a circular orbit,
\[
  I_r=0,
  \qquad
  E=\omega L.
\]
Radial oscillation has twice the angular oscillator frequency: one full radial
cycle occurs while the Cartesian oscillator runs through half of its ellipse.
Therefore the Hamiltonian in polar actions must be
\[
  \boxed{
  H=\omega(2I_r+|I_\theta|).
  }
\]
Solving for the radial action,
\[
  \boxed{
  I_r=\frac{E}{2\omega}-\frac{|L|}{2}.
  }
\]
This formula records the cycle-basis dependence of action variables. Cartesian
actions and polar actions are related by an integral linear change of torus
cycles, not by an arbitrary smooth coordinate change.

\subsection{Symmetric Free Top: Actions And Dense Motion}
\label{sec:detailed-symmetric-top-actions}

For a symmetric free rigid body with
\[
  I_1=I_2>I_3,
\]
use Euler angles \((\phi,\theta,\psi)\). The conserved quantities are
\[
  H,\qquad J^2,\qquad J_z.
\]
The momentum conjugate to \(\psi\) is the body-axis component of angular
momentum:
\[
  p_\psi=M_3.
\]
The momentum conjugate to \(\phi\) is the spatial vertical component:
\[
  p_\phi=J_z.
\]
The total angular momentum satisfies the standard Euler-angle identity
\[
  J^2
  =
  p_\theta^2
  +
  \frac{(p_\phi-p_\psi\cos\theta)^2}{\sin^2\theta}
  +
  p_\psi^2.
\]
This equation is the rigid-body analogue of the angular-momentum identity in
spherical coordinates. It says that \(p_\theta\) is not an independent
constant; it is determined by \(J\), \(J_z=p_\phi\), \(M_3=p_\psi\), and
\(\theta\).

The convention \(I_1=I_2>I_3\) is the prolate convention: the moment about the
symmetry axis is smaller than the transverse moments. An oblate symmetric top
has the reversed inequality. The formulas below adapt by the same algebra with
the corresponding chamber inequalities. For the symmetric top, whose classical
heavy-top version was analyzed by Lagrange in 1788,
\[
  H
  =
  \frac{J^2-M_3^2}{2I_1}
  +
  \frac{M_3^2}{2I_3}.
\]
Solving for \(M_3^2\) gives
\[
  M_3^2
  =
  \frac{2I_1I_3}{I_1-I_3}
  \left(
    H-\frac{J^2}{2I_1}
  \right).
\]
Thus one action is
\[
  I_\psi=M_3
  =
  \sqrt{
  \frac{2I_1I_3}{I_1-I_3}
  \left(
    H-\frac{J^2}{2I_1}
  \right)
  },
\]
with sign chosen by the orientation of the \(\psi\)-cycle. Another action is
\[
  I_\phi=J_z.
\]
The remaining action is the cycle traced by the nutation angle \(\theta\):
\[
  I_\theta
  =
  \frac1{2\pi}\oint p_\theta\,\dd\theta
  =
  \frac1\pi\int_{\theta_-}^{\theta_+}
  \sqrt{
    J^2-M_3^2
    -
    \frac{(J_z-M_3\cos\theta)^2}{\sin^2\theta}
  }\;\dd\theta .
\]
The endpoints \(\theta_\pm\) are the roots where the square root vanishes.
To see the evaluation, set
\[
  a=J_z,\qquad b=M_3,\qquad x=\cos\theta .
\]
Then
\[
  I_\theta
  =
  \frac1\pi\int_{x_-}^{x_+}
  \frac{\sqrt{J^2(1-x^2)-(a-bx)^2}}{1-x^2}\,\dd x,
\]
where
\[
  x_\pm=
  \frac{ab\pm\sqrt{(J^2-a^2)(J^2-b^2)}}{J^2}.
\]
This integral is the symplectic area on the angular-momentum sphere between
the two latitude circles \(J_z=a\) and \(M_3=b\), divided by \(2\pi\). The
cap-area formula is elementary in spherical coordinates. Write the coadjoint
sphere as \(M=Jn\), with
\(n=(\sin\vartheta\cos\varphi,\sin\vartheta\sin\varphi,\cos\vartheta)\).
The Kirillov form is
\[
  \omega_{\mathrm{KKS}}=J\sin\vartheta\,\dd\vartheta\wedge\dd\varphi,
\]
so the cap above \(M_z=c\), where \(\cos\vartheta=c/J\), has area
\[
  \int_0^{2\pi}\int_0^{\arccos(c/J)}
  J\sin\vartheta\,\dd\vartheta\,\dd\varphi
  =
  2\pi(J-c).
\]
Thus a cap of height \(h=J-c\) has symplectic area \(2\pi h\). The two fixed
projections remove caps determined by \(|a|\) and \(|b|\); the remaining
nutation cycle has height \(J-\max\{|a|,|b|\}\). Hence
\[
  \boxed{
  I_\theta=J-\max\{|J_z|,|M_3|\}.
  }
\]
Thus the simple formula
\[
  I_\theta=J-|J_z|
\]
is the local chamber \(|J_z|\ge |M_3|\). In the other chamber the limiting
projection is the body-axis component, and the same cycle convention gives
\(I_\theta=J-|M_3|\). This chamber dependence is not a new physical
conservation law; it records which projection of the angular momentum sphere
sets the polar turning circle. The spin action \(I_\psi=M_3\) remains a
separate action measuring rotation about the body's symmetry axis.
On the chamber wall \(|J_z|=|M_3|\), the two polar circles touch the same
turning latitude and the chosen nutation cycle changes its representative. The
formula is continuous there, but the local action chart changes because the
cycle basis used to describe the torus is changing.

In action variables, this means
\[
  J=I_\theta+\max\{|I_\phi|,|I_\psi|\},
\]
and hence the free symmetric-top Hamiltonian is
\[
  H(I_\theta,I_\phi,I_\psi)
  =
  \frac{
    \big(I_\theta+\max\{|I_\phi|,|I_\psi|\}\big)^2-I_\psi^2
  }{2I_1}
  +
  \frac{I_\psi^2}{2I_3},
\]
within a fixed chamber. The absolute values and maximum are the price paid for
using Euler-angle cycles across coordinate singularities; once a chamber is
chosen, the Hamiltonian is smooth in the corresponding local action variables.

Let \(\alpha_j\) denote the angle variables conjugate to the action variables
\(I_j=(I_\theta,I_\phi,I_\psi)\) in the chosen chamber. The angle equations are
\[
  \dot\alpha_j=\frac{\partial H}{\partial I_j}.
\]
The Hamiltonian does not depend on the inertial direction of the angular momentum. This
degeneracy comes from the continuous symmetry of rotations about the spatial
angular-momentum direction, not from chaotic dynamics. In the nonsymmetric free
rigid body the generic invariant sets are still Liouville tori, but the action
integrals involve elliptic functions in place of the elementary chamber formula
above. Concretely, set
\[
  A=2H,\qquad B=J^2,
\]
and take \(I_1<I_2<I_3\). The two algebraic integrals
\[
  A=I_1\Omega_1^2+I_2\Omega_2^2+I_3\Omega_3^2,\qquad
  B=I_1^2\Omega_1^2+I_2^2\Omega_2^2+I_3^2\Omega_3^2
\]
solve for \(\Omega_1^2\) and \(\Omega_3^2\) in terms of
\(x=\Omega_2^2\):
\[
\begin{aligned}
  \Omega_1^2
  &=
  \frac{B-AI_3-I_2(I_2-I_3)x}{I_1(I_1-I_3)},\\
  \Omega_3^2
  &=
  \frac{AI_1-B-I_2(I_1-I_2)x}{I_3(I_1-I_3)}.
\end{aligned}
\]
Substituting these expressions into
\(I_2\dot\Omega_2=(I_3-I_1)\Omega_3\Omega_1\) gives a quartic polynomial in
\(\Omega_2\). If the two relevant roots in the chosen energy chamber are
ordered as \(a^2\le x\le b^2\), then the equation takes the form
\[
  \dot\Omega_2^2
  =
  \frac{(I_3-I_2)(I_2-I_1)}{I_1I_3}
  (\Omega_2^2-a^2)(b^2-\Omega_2^2)
\]
with
\[
  \{a^2,b^2\}
  =
  \left\{
  \frac{AI_3-B}{I_2(I_3-I_2)},
  \frac{B-AI_1}{I_2(I_2-I_1)}
  \right\},
\]
listed in the order appropriate to the chamber. Hence, after choosing
the origin of time and a chamber in which \(a^2\le \Omega_2^2\le b^2\),
\[
  \Omega_2(t)=b\,\operatorname{dn}(\nu(t-t_0),k),
  \qquad
  \nu=b\sqrt{\frac{(I_3-I_2)(I_2-I_1)}{I_1I_3}},
  \qquad
  k^2=1-\frac{a^2}{b^2},
\]
or an equivalent \(\operatorname{sn}/\operatorname{cn}\) parametrization
depending on the chamber. To check the \(\operatorname{dn}\) form directly,
use
\[
  \frac{\dd}{\dd u}\operatorname{dn}(u,k)
  =
  -k^2\operatorname{sn}(u,k)\operatorname{cn}(u,k),
\]
together with
\[
  b^2\operatorname{dn}^2-a^2=b^2k^2\operatorname{cn}^2,
  \qquad
  b^2-b^2\operatorname{dn}^2=b^2k^2\operatorname{sn}^2.
\]
Then \(\dot\Omega_2^2\) equals the quartic above precisely when
\[
  \nu=b\sqrt{\frac{(I_3-I_2)(I_2-I_1)}{I_1I_3}}.
\]
The elementary symmetric-top formula is therefore a
degenerate limit of the elliptic-function solution.

\section{Associated Demos}
\label{sec:rigid-associated-demos}

The script \path{demos/python/rigid_body_euler_top.py} integrates Euler's
equations, reconstructs the attitude with
\[
  R_{n+1}=R_n\exp(\Delta t\,\widehat\Omega_{n+1/2}),
\]
and reports drift in \(H\), \(C=\norm M^2\), the spatial momentum
\(J=RM\), and the constraints \(R^TR=I\), \(\det R=1\). The Mathematica script
\path{demos/mathematica/RigidBodyEulerTop.wl} gives a symbolic and numerical
companion for the reduced equations. A useful check is to compare the numerical
drift in \(H\), \(C\), and \(J\) for several time steps and initial conditions
near the intermediate axis.
